% Part: lambda-calculus
% Chapter: syntax
% Section: term-revisited

\documentclass[../../../include/open-logic-section]{subfiles}

\begin{document}

\olfileid{lam}{syn}{tr}
\olsection{الحدود بوصفها أصناف تكافؤ بحسب~${\OLId{greek-variable}{alpha}}$}

من الآن ندرس الحدود مع إهمال الفرق الذي لا يتجاوز التكافؤ بحسب~${\OLId{greek-variable}{alpha}}$.
فإذا كتبنا حدًا قصدنا صنف تكافؤه بحسب~${\OLId{greek-variable}{alpha}}$.
فكتابتنا~$\lambd[{\OLId{latin-ordinary}{a}}][\lambd[{\OLId{latin-ordinary}{b}}][{\OLId{latin-ordinary}{a}} {\OLId{latin-ordinary}{c}}]]$ تعني مجموعة الحدود التي تكافئه بحسب~${\OLId{greek-variable}{alpha}}$،
ومنها~$\lambd[{\OLId{latin-ordinary}{a}}][\lambd[{\OLId{latin-ordinary}{b}}][{\OLId{latin-ordinary}{a}} {\OLId{latin-ordinary}{c}}]]$ و$\lambd[{\OLId{latin-ordinary}{b}}][\lambd[{\OLId{latin-ordinary}{a}}][{\OLId{latin-ordinary}{b}} {\OLId{latin-ordinary}{c}}]]$ وما شابهها.

وكانت حروف مثل~${\OLId{latin-looped}{N}}, {\OLId{latin-looped}{Q}}$ في الأقسام الماضية أسماء لحدود مفردة؛
أما فيما يأتي فنسمي بها الأصناف، وهي التي تكون موضوع البحث بدل الحدود المفردة.
وتبقى حروف مثل~${\OLId{latin-ordinary}{x}}, {\OLId{latin-ordinary}{y}}$ أسماء للمتغيرات.

ونجعل~$\rep{{\OLId{latin-looped}{M}}}$ رمزًا لـ!!{element} كيفما اتفق من الصنف~${\OLId{latin-looped}{M}}$.
فإن احتجنا إلى ممثلين متعددين استعملنا~$\rep{{\OLId{latin-looped}{M}}}[{\OLMathNumeral{0}}], \rep{{\OLId{latin-looped}{M}}}[{\OLMathNumeral{1}}]$، وهكذا.

ولإيجاز العبارة نبقي ترميز الحدود، ونعرّف معناه على الأصناف:
\begin{defn}
  \begin{enumerate}
  \item $\lambd[{\OLId{latin-ordinary}{x}}][{\OLId{latin-looped}{N}}]$ هو، بالتعريف، الصنف الذي ينتمي إليه
    $\lambd[{\OLId{latin-ordinary}{x}}][\rep{{\OLId{latin-looped}{N}}}]$.
  \item $PQ$ هو، بالتعريف، الصنف الذي ينتمي إليه~$\rep{{\OLId{latin-looped}{P}}}\rep{{\OLId{latin-looped}{Q}}}$.
  \end{enumerate}
\end{defn}

وهذان التعريفان لا يتوقفان على اختيار الممثل؛ فإن تحويل~${\OLId{greek-variable}{alpha}}$ متوافق، ومن ذلك يظهر حسن تعريفهما.

\begin{defn} \ollabel{def:fv}
  مجموعة \emph{المتغيرات الحرة} لصنف التكافؤ بحسب~${\OLId{greek-variable}{alpha}}$ المسمى~${\OLId{latin-looped}{M}}$،
  ورمزها~$\FV{{\OLId{latin-looped}{M}}}$، هي بالتعريف~$\FV{\rep{{\OLId{latin-looped}{M}}}}$.
\end{defn}

ولا تتعلق هذه المجموعة باختيار الممثل، لأن~$\FV{\rep{{\OLId{latin-looped}{M}}}[{\OLMathNumeral{0}}]} =
\FV{\rep{{\OLId{latin-looped}{M}}}[{\OLMathNumeral{1}}]}$، كما أثبتنا في~\olref[alp]{thm:fv}.

وننقل ترميز الإحلال إلى الأصناف كذلك:
\begin{defn} \ollabel{def:sub}
  \emph{إحلال}~${\OLId{latin-looped}{R}}$ محل~${\OLId{latin-ordinary}{y}}$ في~${\OLId{latin-looped}{M}}$، ورمزه~$\Subst{{\OLId{latin-looped}{M}}}{{\OLId{latin-looped}{R}}}{{\OLId{latin-ordinary}{y}}}$،
  هو صنف التكافؤ بحسب~${\OLId{greek-variable}{alpha}}$ الذي ينتمي إليه
  $\Subst{\rep{{\OLId{latin-looped}{M}}}}{\rep{{\OLId{latin-looped}{R}}}}{{\OLId{latin-ordinary}{y}}}$؛ ونأخذ~$\rep{{\OLId{latin-looped}{M}}}$ و$\rep{{\OLId{latin-looped}{R}}}$ أي ممثلين يكون الإحلال عليهما معرّفًا.
\end{defn}

وحسن تعريف هذا الإحلال ثابت أيضًا بـ\olref[alp]{cor:sub}.

ويتضح مقدار ما ييسره هذا التعريف من الاستدلال في المثال الآتي:
\begin{align}
  \Subst{\lambd[{\OLId{latin-ordinary}{x}}][{\OLId{latin-ordinary}{x}}]}{{\OLId{latin-ordinary}{y}}}{{\OLId{latin-ordinary}{x}}} & \ollabel{eq:1}\\
  &= \Subst{\lambd[{\OLId{latin-ordinary}{z}}][{\OLId{latin-ordinary}{z}}]}{{\OLId{latin-ordinary}{y}}}{{\OLId{latin-ordinary}{x}}} \ollabel{eq:2}\\
  &= \lambd[{\OLId{latin-ordinary}{z}}][\Subst{{\OLId{latin-ordinary}{z}}}{{\OLId{latin-ordinary}{y}}}{{\OLId{latin-ordinary}{x}}}] \\
  &= \lambd[{\OLId{latin-ordinary}{z}}][{\OLId{latin-ordinary}{z}}]
\end{align}

لو حملنا~\olref{eq:1} على الإحلال في الحدود المفردة لكان غير معرّف.
لكن المراد الآن الإحلال في الأصناف، كما تقدم؛ ولذلك جاز الانتقال إلى~\olref{eq:2}:
فلنا أن نستبدل~$\lambd[{\OLId{latin-ordinary}{x}}][{\OLId{latin-ordinary}{x}}]$ بـ$\lambd[{\OLId{latin-ordinary}{z}}][{\OLId{latin-ordinary}{z}}]$، إذ كلاهما في صنف واحد.

وللعلة نفسها نفترض فيما يأتي أن اختيار الممثلين قد استوفى شروط الإحلال.
فإذا ورد~$\Subst{\lambd[{\OLId{latin-ordinary}{x}}][{\OLId{latin-looped}{N}}]}{{\OLId{latin-looped}{R}}}{{\OLId{latin-ordinary}{y}}}$، مثلًا، حملناه على أن الممثل~$\lambd[{\OLId{latin-ordinary}{x}}][{\OLId{latin-looped}{N}}]$
اختير بحيث~${\OLId{latin-ordinary}{x}} \neq {\OLId{latin-ordinary}{y}}$ و${\OLId{latin-ordinary}{x}} \notin \FV{{\OLId{latin-looped}{R}}}$.

ولما كان إطلاق «صنف» على~$\lambd[{\OLId{latin-ordinary}{x}}][{\OLId{latin-ordinary}{x}}]$ غير مألوف بعض الشيء،
سمينا هذه الأشياء فيما يأتي حدود~$\Lambda$، أو «حدودًا» بإطلاق في بقية هذا الجزء،
تمييزًا لها من حدود~$\lambda$ المعهودة.

\begin{editorial}
  لا نستغني بذلك عن الحدود الأصلية؛ فتعريف حدود~$\Lambda$ كله يرجع إلى
  حدود~$\lambda$. ولم نضع بعد طريقة لتعريف الدوال على حدود~$\Lambda$،
  فلا بد من تعريف كل دالة أولًا على حدود~$\lambda$،
  ثم إنزالها إلى حدود~$\Lambda$، على نحو ما صنعنا بالإحلال.
  غير أننا نفترض أن القارئ يدرك حدسيًا كيف تعرّف الدوال على حدود~$\Lambda$.
\end{editorial}
\end{document}
